\chapter{Brownian Motion}

\begin{center}\itshape
In 1827 Robert Brown watched particles from pollen grains jitter in water.\\
In 1905 Einstein explained the jitter by molecular impacts,\\
and in 1923 Wiener built the curve such a particle follows:\\
continuous everywhere, differentiable nowhere.
\end{center}

Shrink the steps of a random walk to $\frac1{\sqrt n}$ and speed up time by $n$: by the central limit theorem, the position at time $t$
becomes approximately $N(0,t)$, the displacements over disjoint time intervals stay independent, and the path looks
more and more like a continuous curve. The limiting object is \emph{Brownian motion}, the most important random
process in mathematics. It models the motion of a particle in a fluid, the noise in an electrical circuit, and, since
Louis Bachelier's thesis of 1900, the fluctuations of stock prices.

The first question is whether it exists at all: can one have continuous paths with independent Gaussian increments?
We answer with Paul L\'evy's construction, which refines a path by random midpoint displacements, and Borel--Cantelli
shows that the refinements converge uniformly. Then we study the paths. Their scaling makes them self-similar; their
quadratic variation over $[0,t]$ is exactly $t$; they have no derivative at any point. The reflection principle of Chapter~7
passes to the limit and gives the law of the maximum. A discrete version of It\^o's formula shows where the extra
term of stochastic calculus comes from. Finally we price an option in the binomial model of Cox, Ross and
Rubinstein and let the time step go to zero: the central limit theorem turns the binomial price into the formula of
Black and Scholes.

\begin{think}
Zoom into a Brownian path, magnifying time by a factor $4$. By what factor must you magnify space for the picture to
look statistically the same? What does your answer say about the slope of the path?
\end{think}

\section{From random walks to Brownian motion}

Let $S_k$ be a random walk with independent steps of mean $0$ and variance $1$, and set $W^{(n)}_t=\frac{S_{\lfloor nt\rfloor}}{\sqrt n}$. For $s<t$, the increment
$W^{(n)}_t-W^{(n)}_s$ is a sum of $\lfloor nt\rfloor-\lfloor ns\rfloor$ steps divided by $\sqrt n$, so by the central limit theorem it converges in law to $N(0,t-s)$, and
increments over disjoint intervals are independent. This suggests the following definition.

\begin{definition}\label{def:bm}
A \emph{Brownian motion} is a family of random variables $(B_t)_{t\ge0}$ on one probability space such that
\begin{enumerate}[label=(\roman*)]
\item $B_0=0$, and for every $\omega$ the path $t\mapsto B_t(\omega)$ is continuous;
\item for $0\le t_0<t_1<\dots<t_k$, the increments $B_{t_1}-B_{t_0},\dots,B_{t_k}-B_{t_{k-1}}$ are independent;
\item for $s<t$, $B_t-B_s$ is $N(0,t-s)$.
\end{enumerate}
\end{definition}

\begin{proposition}\label{prop:bmgauss}
Let $(B_t)$ have continuous paths with $B_0=0$. Then (ii) and (iii) hold if and only if for all $t_1,\dots,t_k$ the vector $(B_{t_1},\dots,B_{t_k})$ is a centred Gaussian
vector with $\operatorname{Cov}(B_s,B_t)=\min(s,t)$.
\end{proposition}
\begin{proof}
If (ii) and (iii) hold, the increments over $0<t_1<\dots<t_k$ form a Gaussian vector with independent components (Theorem~\ref{thm:gaussvec}(b),(c)), and $(B_{t_i})$
is a linear image of it, hence Gaussian (Theorem~\ref{thm:gaussvec}(d)); for $s<t$, $\operatorname{Cov}(B_s,B_t)=\operatorname{Var}(B_s)+\operatorname{Cov}(B_s,B_t-B_s)=s$. Conversely, the increments are a
linear image of $(B_{t_i})$, hence a Gaussian vector, and for $s<t\le u<v$,
\[
\operatorname{Cov}(B_t-B_s,B_v-B_u)=t-s-t+s=0,
\]
while $\operatorname{Var}(B_t-B_s)=t-2s+s=t-s$. Theorem~\ref{thm:gaussvec}(c) gives (ii), and (iii) follows.
\end{proof}

That the rescaled walks converge to Brownian motion as whole paths, and not only through finitely many times, is
\emph{Donsker's invariance principle} (1951), proved in Tome~10. Here we need a construction: nothing so far guarantees that
continuous paths with these increments exist.

\section{L\'evy's construction}

Let $D_n=\{k2^{-n}:0\le k\le2^n\}$ be the dyadic points of level $n$ in $[0,1]$, and $D=\bigcup_nD_n$. Take independent $N(0,1)$ variables $Z_t$, one for each $t\in D$
(admitted fact (A1) of Chapter~5). Define $B$ on $D$ level by level: $B_0=0$, $B_1=Z_1$, and for $t\in D_n\setminus D_{n-1}$, $n\ge1$,
\[
B_t=\frac{B_{t-2^{-n}}+B_{t+2^{-n}}}2+2^{-(n+1)/2}Z_t.
\]
Each new point is placed at the midpoint of its two neighbours, then displaced by an independent Gaussian of the right
size (Figure~\ref{fig:levy}).

\begin{lemma}\label{lem:levylevel}
For every $n$, the $2^n$ increments $B_{k2^{-n}}-B_{(k-1)2^{-n}}$ are independent $N(0,2^{-n})$ variables.
\end{lemma}
\begin{proof}
By induction; $n=0$ is clear. All the variables involved are linear combinations of the independent Gaussians $Z_t$, so the
increments of level $n$ form a Gaussian vector, and it suffices to compute variances and covariances
(Theorem~\ref{thm:gaussvec}(c),(d)). An interval of level $n-1$ with increment $\Delta$ (of variance $2^{-n+1}$) is split at its midpoint $t$ into two halves with
increments $\frac\Delta2\pm2^{-(n+1)/2}Z_t$. Their variances are $\frac{2^{-n+1}}4+2^{-n-1}=2^{-n}$, and their covariance is $\frac{2^{-n+1}}4-2^{-n-1}=0$. Halves of different intervals of
level $n-1$ are built from independent variables (the level-$(n-1)$ increments, independent by induction, and distinct $Z_t$), so they are
uncorrelated too.
\end{proof}

\begin{theorem}[Wiener; L\'evy's construction]\label{thm:levy}
Brownian motion exists.
\end{theorem}
\begin{proof}
\emph{Uniform convergence.} Let $F_n$ be the function on $[0,1]$ that equals $B$ on $D_n$ and is linear between consecutive points of $D_n$. Then $F_n-F_{n-1}$
is piecewise linear, vanishes on $D_{n-1}$, and equals $2^{-(n+1)/2}Z_t$ at each $t\in D_n\setminus D_{n-1}$, so
\[
\|F_n-F_{n-1}\|_\infty=2^{-(n+1)/2}\max_{t\in D_n\setminus D_{n-1}}|Z_t|.
\]
For $x\ge1$, $P(|Z|\ge x)=2\int_x^\infty\frac{e^{-u^2/2}}{\sqrt{2\pi}}\dd u\le2\int_x^\infty\frac ux\frac{e^{-u^2/2}}{\sqrt{2\pi}}\dd u=\frac{2e^{-x^2/2}}{x\sqrt{2\pi}}\le e^{-x^2/2}$. Hence
\[
P\Bigl(\max_{t\in D_n\setminus D_{n-1}}|Z_t|\ge2\sqrt n\Bigr)\le2^{n-1}e^{-2n},
\]
which is summable. By Borel--Cantelli (Theorem~\ref{thm:bc}), almost surely $\max|Z_t|<2\sqrt n$ for all large $n$, and then $\sum_n\|F_n-F_{n-1}\|_\infty\le\sum_n2^{-(n+1)/2}2\sqrt n<\infty$.
So $F_n$ converges uniformly on $[0,1]$ to a continuous function, which equals $B$ on $D$ (as $F_m=B$ on $D_n$ for $m\ge n$). On the null event where
this fails, set the limit to $0$. Call the limit $B$ again.

\emph{The law.} Let $t_1,\dots,t_k\in[0,1]$ and $t_i^{(n)}=\lceil2^nt_i\rceil2^{-n}\in D_n$. By Lemma~\ref{lem:levylevel}, as in the proof of Proposition~\ref{prop:bmgauss},
$(B_{t_i^{(n)}})_i$ is a centred Gaussian vector with covariances $\min(t^{(n)}_i,t^{(n)}_j)$, so its characteristic function is $e^{-u^{\mathsf T}\Sigma_nu/2}$. As $n\to\infty$, $B_{t_i^{(n)}}\to B_{t_i}$ by
continuity, and by dominated convergence the characteristic function of $(B_{t_i})$ is $e^{-u^{\mathsf T}\Sigma u/2}$ with $\Sigma_{ij}=\min(t_i,t_j)$. By Theorem~\ref{thm:inversion}(a)
and Theorem~\ref{thm:gaussvec}(a),(b), $(B_{t_i})$ is a centred Gaussian vector with these covariances, so $B$ is a Brownian motion on $[0,1]$
(Proposition~\ref{prop:bmgauss}).

\emph{All times.} Take independent copies $B^{(0)},B^{(1)},\ldots$ of this process on $[0,1]$ (built from disjoint families of $Z$'s), and for $t\in[j,j+1]$ set
$B_t=B^{(0)}_1+\dots+B^{(j-1)}_1+B^{(j)}_{t-j}$. The paths are continuous, and an increment over $[s,t]$ is a sum of increments of independent copies over the
pieces of $[s,t]$ cut at the integers: independent Gaussians with variances adding up to $t-s$ (Proposition~\ref{prop:gausssum}). Increments over disjoint
intervals are built from disjoint pieces, hence independent.
\end{proof}

\begin{figure}[ht]
\centering
\fig{ch11-levy}
\caption{L\'evy's construction: the piecewise linear paths $F_0$, $F_2$, $F_4$ and $F_{10}$ (darkest), each obtained from the previous one by random
midpoint displacements of size $2^{-(n+1)/2}$.}
\label{fig:levy}
\end{figure}

\begin{beauty}
The construction is also a fast algorithm: a path on $2^n$ points costs $2^n$ Gaussian draws, and can be refined later
where more detail is needed, without changing what is already drawn. The same midpoint displacement, with other
scaling exponents, draws the fractal mountains of computer graphics.
\end{beauty}

\section{Scaling, symmetry and the Markov property}

Let $\mathcal F_s=\sigma(B_u:u\le s)$ be the information revealed by the path up to time $s$.

\begin{proposition}\label{prop:scaling}
Let $B$ be a Brownian motion, $c>0$ and $s\ge0$. Then the following are Brownian motions:
\begin{enumerate}[label=(\alph*)]
\item (\emph{scaling}) $\frac1cB_{c^2t}$;
\item (\emph{symmetry}) $-B_t$;
\item (\emph{Markov property}) $B_{s+t}-B_s$, and this process is independent of $\mathcal F_s$.
\end{enumerate}
\end{proposition}
\begin{proof}
In each case the paths are continuous and start at $0$, and the times $t_1,\dots,t_k$ give a centred Gaussian vector (a linear image of
one, Theorem~\ref{thm:gaussvec}(d)). By Proposition~\ref{prop:bmgauss} it suffices to compute covariances: $\frac1{c^2}\min(c^2s,c^2t)=\min(s,t)$; $\min(s,t)$;
and $\min(s+t,s+t')-s-s+s=\min(t,t')$. For the independence in (c), take $u_1,\dots,u_j\le s$ and $t_1,\dots,t_k\ge0$. The vector made of the $B_{u_i}$ and the
$B_{s+t_l}-B_s$ is Gaussian, and $\operatorname{Cov}(B_{u_i},B_{s+t_l}-B_s)=u_i-u_i=0$; so its characteristic function factorises (Theorem~\ref{thm:gaussvec}(a)), and the two
blocks are independent (Theorem~\ref{thm:inversion}(a)). Events determined by finitely many times form $\pi$-systems that generate the two
$\sigma$-algebras, and Theorem~\ref{thm:uniqueness} extends the independence to them.
\end{proof}

Scaling answers the opening question: magnifying time by $4$ requires magnifying space by $2$. A Brownian path is
statistically self-similar (Figure~\ref{fig:zoom}). Over a time $h$ it moves by about $\sqrt h$, so its difference quotients are of order $\frac1{\sqrt h}$, which
explodes: we shall see that the path has no derivative anywhere.

\begin{figure}[ht]
\centering
\fig{ch11-zoom}
\caption{One Brownian path on $[0,1]$ (left), on $[0,\frac14]$ magnified $2$ times in space (middle), and on $[0,\frac1{16}]$ magnified $4$ times
(right). The three pictures have the same law.}
\label{fig:zoom}
\end{figure}

\begin{lemma}\label{lem:bmmax}
For $\lambda>0$, $P\bigl(\max_{t\le1}|B_t|>\lambda\bigr)\le\frac1{\lambda^2}$.
\end{lemma}
\begin{proof}
For fixed $m$, $(B_{k2^{-m}})_{0\le k\le2^m}$ is a martingale, a sum of independent centred Gaussians, with $\mathbb E[B_1^2]=1$; by Kolmogorov's inequality
(Corollary~\ref{cor:kolmogorov}), $P(\max_{t\in D_m}|B_t|\ge\lambda)\le\frac1{\lambda^2}$. These events increase with $m$, and by continuity their union contains
$\{\max_{t\le1}|B_t|>\lambda\}$, as the maximum is approached along dyadic times.
\end{proof}

\begin{proposition}[time inversion]\label{prop:inversion}
$\frac{B_t}t\to0$ almost surely as $t\to\infty$, and the process $X_0=0$, $X_t=tB_{1/t}$ ($t>0$) is a Brownian motion.
\end{proposition}
\begin{proof}
$B_n$ is a sum of $n$ independent $N(0,1)$ variables, so $\frac{B_n}n\to0$ almost surely by the strong law (Theorem~\ref{thm:slln4}). By Proposition~\ref{prop:scaling}(c) and
Lemma~\ref{lem:bmmax}, $P(\max_{t\in[n,n+1]}|B_t-B_n|>n^{3/4})\le n^{-3/2}$, which is summable, so by Borel--Cantelli $\max_{[n,n+1]}|B_t-B_n|\le n^{3/4}$ for all large $n$. So for $t\in[n,n+1]$,
$\frac{|B_t|}t\le\frac{|B_n|+n^{3/4}}n\to0$.

The process $X$ is continuous on $]0,\infty[$, and at $0$ since $X_t=\frac{B_u}u$ with $u=\frac1t\to\infty$. Its values at finitely many times form a centred Gaussian
vector (a linear image of one) with $\operatorname{Cov}(X_s,X_t)=st\min(\frac1s,\frac1t)=\min(s,t)$; conclude by Proposition~\ref{prop:bmgauss}.
\end{proof}

\section{Quadratic variation and roughness}

For a partition $\pi:0=t_0<t_1<\dots<t_m=t$ of $[0,t]$, with mesh $|\pi|=\max_k(t_k-t_{k-1})$, let
\[
Q(\pi)=\sum_{k=1}^m\bigl(B_{t_k}-B_{t_{k-1}}\bigr)^2.
\]
For a function with a continuous derivative, $Q(\pi)\le\max|f'|^2\,|\pi|\,t\to0$. Not so for Brownian motion.

\begin{theorem}[quadratic variation]\label{thm:qv}
$\mathbb E\bigl[(Q(\pi)-t)^2\bigr]=2\sum_k(t_k-t_{k-1})^2\le2t|\pi|$, so $Q(\pi)\to t$ in $L^2$ as $|\pi|\to0$. For the dyadic partitions $\pi_n$ of mesh $t2^{-n}$, $Q(\pi_n)\to t$
almost surely.
\end{theorem}
\begin{proof}
The terms $Y_k=(B_{t_k}-B_{t_{k-1}})^2-(t_k-t_{k-1})$ are independent, centred, and $\operatorname{Var}(Y_k)=2(t_k-t_{k-1})^2$, since for $N$ of law $N(0,h)$, $\mathbb E[N^4]=3h^2$. So
$\mathbb E[(Q(\pi)-t)^2]=\operatorname{Var}(\sum Y_k)=2\sum(t_k-t_{k-1})^2\le2|\pi|\sum(t_k-t_{k-1})$. For $\pi_n$ this is $2t^22^{-n}$; by Chebyshev, $P(|Q(\pi_n)-t|>\varepsilon)\le\frac{2t^2}{\varepsilon^22^n}$,
summable, and Borel--Cantelli gives the almost sure convergence.
\end{proof}

\begin{corollary}\label{cor:variation}
Almost surely, the path has infinite total variation on every interval $[s,t]$ with $s<t$: the lengths $\sum_k|B_{t_k}-B_{t_{k-1}}|$ over partitions of $[s,t]$ are
unbounded.
\end{corollary}
\begin{proof}
By Proposition~\ref{prop:scaling}(c), Theorem~\ref{thm:qv} applies on $[s,t]$; almost surely the dyadic sums converge to $t-s>0$ for all rational $s<t$ at once. If
the variation $V$ on $[s,t]$ were finite, then $Q(\pi_n)\le\max_k|B_{t_k}-B_{t_{k-1}}|\cdot V\to0$ by uniform continuity, a contradiction. Every interval contains one
with rational endpoints.
\end{proof}

So $\int f\dd B$ cannot be defined path by path as a Stieltjes integral. Stochastic calculus (Tome~11) defines it in $L^2$,
using precisely the rule ``$(\dd B)^2=\dd t$'' expressed by Theorem~\ref{thm:qv}.

\begin{theorem}[Paley, Wiener and Zygmund, 1933]\label{thm:pwz}
Almost surely, the Brownian path is differentiable at no point of $[0,\infty[$.
\end{theorem}
\begin{proof}
By Proposition~\ref{prop:scaling}(c) it suffices to treat $s\in[0,1]$ (then shift by integers). Suppose the path is differentiable at some $s\in[0,1]$. Then the
difference quotients are bounded near $s$, and by continuity on $[0,2]$ there is an integer $M$ with $|B_t-B_s|\le M|t-s|$ for all $t\in[0,2]$. For $n\ge2$ let
$k=\lfloor2^ns\rfloor$, so that $k2^{-n}\le s<(k+1)2^{-n}$. For $j=1,2,3$, the points $(k+j-1)2^{-n}$ and $(k+j)2^{-n}$ lie within $(j+1)2^{-n}$ of $s$, so
\[
\begin{aligned}\bigl|B_{(k+j)2^{-n}}-B_{(k+j-1)2^{-n}}\bigr|&\le M(2j+1)2^{-n}\eqbr\le7M2^{-n}.\end{aligned}
\]
Let $A_{M,n}$ be the event that for some $k\in\{0,\dots,2^n\}$ three consecutive increments of length $2^{-n}$ starting at $k2^{-n}$ are all at most $7M2^{-n}$ in absolute value.
The increments are independent $N(0,2^{-n})$, and $P(|N(0,2^{-n})|\le x)\le\frac{2x}{\sqrt{2\pi}\,2^{-n/2}}\le2^{n/2}x$ (the density is at most $\frac{2^{n/2}}{\sqrt{2\pi}}$), so
\[
\begin{aligned}P(A_{M,n})&\le(2^n+1)\bigl(2^{n/2}\cdot7M2^{-n}\bigr)^3\eqbr\le2(7M)^32^{-n/2},\end{aligned}
\]
which tends to $0$. We showed that differentiability at some $s\in[0,1]$ implies $A_{M,n}$ for some $M$ and all $n\ge2$. But $P(\bigcap_{n\ge2}A_{M,n})\le\inf_nP(A_{M,n})=0$ for each $M$, and a
countable union of null events is null.
\end{proof}

The short proof is due to Dvoretzky, Erd\H os and Kakutani (1961). Weierstrass had astonished his contemporaries in 1872 with
one explicit continuous nowhere differentiable function; the theorem says that, in the sense of Wiener's measure on
paths, almost every continuous function is of this kind.

\section{Martingales and the reflection principle}

A family $(X_t)_{t\ge0}$ adapted to $(\mathcal F_t)$, with $X_t$ integrable, is a martingale if $\mathbb E[X_t\mid\mathcal F_s]=X_s$ for $s\le t$, as in Chapter~10.

\begin{proposition}\label{prop:bmmart}
$B_t$, $B_t^2-t$ and $e^{\theta B_t-\theta^2t/2}$ (for real $\theta$) are martingales.
\end{proposition}
\begin{proof}
Write $B_t=B_s+N$ with $N=B_t-B_s$ of law $N(0,t-s)$, independent of $\mathcal F_s$ (Proposition~\ref{prop:scaling}(c)). By the rules of Proposition~\ref{prop:cerules},
$\mathbb E[B_t\mid\mathcal F_s]=B_s$; $\mathbb E[B_t^2\mid\mathcal F_s]=B_s^2+2B_s\mathbb E[N]+\mathbb E[N^2]=B_s^2+t-s$; and $\mathbb E[e^{\theta B_t}\mid\mathcal F_s]=e^{\theta B_s}\mathbb E[e^{\theta N}]=e^{\theta B_s+\theta^2(t-s)/2}$, using the
Gaussian moment generating function $\mathbb E[e^{\theta N}]=e^{\theta^2(t-s)/2}$ (complete the square).
\end{proof}

Chapter~7 proved that the maximum of a simple random walk is distributed like $|S_n|$, roughly, by reflecting paths after
their first visit to a level. Here is the continuous version, promised there.

\begin{theorem}[reflection principle]\label{thm:reflection}
Let $M_t=\max_{s\le t}B_s$. For $a>0$,
\[
P(M_t\ge a)=2P(B_t\ge a)=P(|B_t|\ge a).
\]
So $M_t$ has the law of $|B_t|$.
\end{theorem}
\begin{proof}
By scaling we may take $t=1$. Let $T=\inf\{s:B_s=a\}$, the hitting time of $a$ ($\infty$ if $a$ is never reached). By continuity $M_1\ge a$ if and only if
$T\le1$, and for every $s$, $\{T\le s\}=\{\max_{u\le s}B_u\ge a\}$ is in $\mathcal F_s$ (the maximum over $[0,s]$ is a supremum over rational times).

\emph{A dyadic version of $T$.} Let $T_n=\lceil2^nT\rceil2^{-n}$, the first dyadic time of level $n$ at or after $T$. The event $\{T_n=k2^{-n}\}=\{(k-1)2^{-n}<T\le k2^{-n}\}$ is in
$\mathcal F_{k2^{-n}}$. For $k2^{-n}<1$, the increment $B_1-B_{k2^{-n}}$ is independent of $\mathcal F_{k2^{-n}}$ and has a symmetric law (Proposition~\ref{prop:scaling}(c)). So for
$\varepsilon>0$,
\[
\begin{aligned}&P\bigl(T_n=k2^{-n},\ B_1-B_{k2^{-n}}>\varepsilon\bigr)\eqbr=P\bigl(T_n=k2^{-n},\ B_1-B_{k2^{-n}}<-\varepsilon\bigr),\end{aligned}
\]
and summing over $k$: $P(T_n<1,\ B_1-B_{T_n}>\varepsilon)=P(T_n<1,\ B_1-B_{T_n}<-\varepsilon)$.

\emph{Passing to the limit.} As $n\to\infty$, $T_n\downarrow T$, so $\mathbf 1_{\{T_n<1\}}\to\mathbf 1_{\{T<1\}}$, and on $\{T<1\}$, $B_{T_n}\to B_T=a$ by continuity. The indicator functions converge
except on $\{B_1=a\pm\varepsilon\}$, a null event, so dominated convergence gives
\[
P(T<1,\ B_1>a+\varepsilon)=P(T<1,\ B_1<a-\varepsilon).
\]
Let $\varepsilon\downarrow0$: $P(T<1,\ B_1>a)=P(T<1,\ B_1<a)$. Now $\{B_1>a\}\subseteq\{T<1\}$ by the intermediate value theorem, and $P(T=1)\le P(B_1=a)=0$. Hence, as $P(B_1=a)=0$,
\[
\begin{aligned}P(M_1\ge a)&=P(T<1)\eqbr=P(B_1>a)+P(T<1,\ B_1<a)\eqbr=2P(B_1>a),\end{aligned}
\]
and $2P(B_1\ge a)=P(|B_1|\ge a)$ by symmetry.
\end{proof}

\begin{figure}[ht]
\centering
\fig{ch11-reflect}
\caption{The reflection principle: after the first time $T$ the path (blue) reaches the level $a$, its reflection in that level
(orange) is another Brownian path. Paths that reach $a$ and end below it correspond to paths that end above it.}
\label{fig:bmreflect}
\end{figure}

\begin{corollary}\label{cor:hitting}
The hitting time $T_a$ of $a>0$ is finite almost surely, with $P(T_a\le t)=2\bigl(1-\Phi(\frac a{\sqrt t})\bigr)$, and $\mathbb E[T_a]=\infty$.
\end{corollary}
\begin{proof}
$P(T_a\le t)=P(M_t\ge a)=2P(B_t\ge a)=2(1-\Phi(\frac a{\sqrt t}))$, which tends to $1$ as $t\to\infty$. And $P(T_a>t)=P(|B_1|<\frac a{\sqrt t})\ge\frac{2a}{\sqrt t}\cdot\frac{e^{-1/2}}{\sqrt{2\pi}}$ for $t\ge a^2$
(the density of $|B_1|$ is at least $\frac{2e^{-1/2}}{\sqrt{2\pi}}$ on $[0,1]$), so $\mathbb E[T_a]=\int_0^\infty P(T_a>t)\dd t=\infty$.
\end{proof}

As for the simple random walk (Corollary~\ref{cor:return1d}), Brownian motion reaches every level, but takes an infinite expected time to do so.

\section{A first look at It\^o's formula}

If $f$ is smooth and $x(t)$ a smooth path, $f(x(t))-f(x(0))=\int_0^tf'(x(s))x'(s)\dd s$. For a random walk, the chain rule acquires a second term.

\begin{proposition}[discrete It\^o formula]\label{prop:discreteito}
Let $S_n=X_1+\dots+X_n$ be the simple random walk, $f:\Z\to\R$, $Df(x)=\frac{f(x+1)-f(x-1)}2$ and $\Delta f(x)=f(x+1)-2f(x)+f(x-1)$. Then
\[
\begin{aligned}f(S_n)&=f(0)+\sum_{k=1}^nDf(S_{k-1})X_k\eqbr+\frac12\sum_{k=1}^n\Delta f(S_{k-1}).\end{aligned}
\]
The first sum is a martingale, so $f(S_n)-\frac12\sum_{k\le n}\Delta f(S_{k-1})$ is a martingale.
\end{proposition}
\begin{proof}
It suffices to check one step: with $x=S_{k-1}$, $f(x+X_k)-f(x)=Df(x)X_k+\frac12\Delta f(x)$. For $X_k=1$ the right side is
$\frac{f(x+1)-f(x-1)+f(x+1)-2f(x)+f(x-1)}2=f(x+1)-f(x)$, and for $X_k=-1$ it is $f(x-1)-f(x)$ in the same way. The first sum is a martingale transform
of the walk by the predictable sequence $Df(S_{k-1})$, which up to time $n$ is bounded by $\max_{|x|<n}|Df(x)|$ (Proposition~\ref{prop:transform}).
\end{proof}

With $f(x)=x^2$, $\Delta f=2$ and we recover the martingale $S_n^2-n$; with $f(x)=x^3$, $\Delta f(x)=6x$, so $S_n^3-3\sum_{k<n}S_k$ is a martingale.

\begin{remark}
Scale space by $\sqrt h$ and time by $h$ and let $h\to0$. The sum $\sum Df(S_{k-1})X_k$ becomes a stochastic integral $\int_0^tf'(B_s)\dd B_s$, and the second
sum, where $\Delta f\approx h f''$ and there are $\frac th$ terms, becomes $\frac12\int_0^tf''(B_s)\dd s$. This is \emph{It\^o's formula},
\[
f(B_t)=f(0)+\int_0^tf'(B_s)\dd B_s+\frac12\int_0^tf''(B_s)\dd s,
\]
proved in Tome~11. The correction term is the trace of the quadratic variation: the path is so rough that the second order
of Taylor's formula survives. In the same way, the discrete Dirichlet problem of Theorem~\ref{thm:dirichlet} becomes the continuous one: a
function harmonic in a domain is the expected value of its boundary values at the exit point of Brownian motion
(Kakutani, 1944), also in Tome~11.
\end{remark}

\section{The binomial model and the Black--Scholes formula}

A bank account grows at a continuous rate $r$: one euro becomes $e^{rt}$ euros at time $t$. A stock has price $S_t$. A \emph{European
put} with strike $K$ and maturity $T$ pays $(K-S_T)^+$ at time $T$ (the right to sell the stock at the price $K$); a \emph{call} pays $(S_T-K)^+$.
What is the fair price of such a contract today?

In the \emph{binomial model} of Cox, Ross and Rubinstein, time is cut into $n$ periods of length $h=\frac Tn$, and in each period the
price is multiplied by $u=e^{\sigma\sqrt h}$ or by $d=e^{-\sigma\sqrt h}$, where $\sigma>0$ is the \emph{volatility}. We assume $d<e^{rh}<u$, which holds as soon as $h$ is small: otherwise
one of the two assets would dominate the other. A \emph{strategy} holds, during each period, a number of shares and an amount
in the bank chosen at the start of the period (a predictable choice, as in Chapter~10), and is \emph{self-financing} if it never
adds or withdraws money. A contract is \emph{replicated} by a strategy whose final value equals its payoff in every
scenario. Its fair price is then the initial cost of the strategy: any other price would allow a riskless profit
(an \emph{arbitrage}), by selling the dearer of the two and buying the cheaper.

\begin{proposition}[one period]\label{prop:oneperiod}
Let $p^*=\frac{e^{rh}-d}{u-d}\in\,]0,1[$. A contract paying $V_u$ if the price goes up and $V_d$ if it goes down, after one period from price $x$, is replicated by
holding $\delta=\frac{V_u-V_d}{x(u-d)}$ shares and the rest in the bank, at the initial cost
\[
V=e^{-rh}\bigl(p^*V_u+(1-p^*)V_d\bigr).
\]
\end{proposition}
\begin{proof}
With $\delta$ shares and $\beta$ in the bank, the final value is $\delta xu+\beta e^{rh}$ or $\delta xd+\beta e^{rh}$. Equating with $V_u$ and $V_d$ gives $\delta$ as stated and
$\beta=e^{-rh}\frac{uV_d-dV_u}{u-d}$. The cost is
\[
\begin{aligned}\delta x+\beta&=\frac{e^{-rh}}{u-d}\bigl((e^{rh}-d)V_u+(u-e^{rh})V_d\bigr)\eqbr=V.\end{aligned}
\]
\end{proof}

\begin{theorem}[binomial pricing]\label{thm:binomial}
In the $n$-period model, a contract paying $g(S_T)$ is replicated by a self-financing strategy with initial cost
\[
V_0=e^{-rT}\,\mathbb E^*\bigl[g(S_T)\bigr],
\]
where under the probability $P^*$ the moves are independent, up with probability $p^*$. Under $P^*$ the discounted price $e^{-rkh}S_{kh}$ is a martingale.
\end{theorem}
\begin{proof}
Define $V_n(x)=g(x)$ and, backwards,
\[
V_k(x)=e^{-rh}\bigl(p^*V_{k+1}(xu)+(1-p^*)V_{k+1}(xd)\bigr).
\]
At time $k$, with price $x$, Proposition~\ref{prop:oneperiod} gives a portfolio of value $V_k(x)$
that becomes $V_{k+1}(S_{(k+1)h})$ at the end of the period, whatever happens; chaining these portfolios is a self-financing strategy from
$V_0(S_0)$ to $g(S_T)$. Unrolling the recursion, $V_k(x)=e^{-r(n-k)h}\mathbb E^*[g(xu^Jd^{n-k-J})]$ with $J$ binomial$(n-k,p^*)$, and $k=0$ gives the formula. Finally
$p^*u+(1-p^*)d=e^{rh}$, so $\mathbb E^*[S_{(k+1)h}\mid\mathcal F_k]=e^{rh}S_{kh}$.
\end{proof}

\begin{figure}[ht]
\centering
\fig{ch11-tree}
\caption{The binomial tree for three periods: from each node the price is multiplied by $u$ or by $d=\frac1u$. Pricing runs
backwards from the payoffs at the leaves, each node receiving the discounted $p^*$-average of its two children.}
\label{fig:tree}
\end{figure}

\begin{beauty}
The real probability that the stock goes up appears nowhere. The price is an expectation, but under the artificial
probability $P^*$ that makes the stock a fair game after discounting. Two traders who disagree completely about the
future of the stock must still agree on the price of the option, because the option can be manufactured from the
stock and the bank.
\end{beauty}

Let $\Phi$ be the $N(0,1)$ distribution function and
\[
d_\pm=\frac{\ln(S_0/K)+(r\pm\frac12\sigma^2)T}{\sigma\sqrt T}.
\]

\begin{theorem}[Black--Scholes as a limit]\label{thm:blackscholes}
As $n\to\infty$, the binomial prices of the put and of the call converge to
\[
\begin{gathered}P_{\mathrm{BS}}=Ke^{-rT}\Phi(-d_-)-S_0\Phi(-d_+),\qq C_{\mathrm{BS}}=S_0\Phi(d_+)-Ke^{-rT}\Phi(d_-).\end{gathered}
\]
\end{theorem}
\begin{proof}
\emph{The law of $\ln S_T$.} Write $a=\sigma\sqrt h$, so $a^2=\sigma^2h$, and $\ln\frac{S_T}{S_0}=a\sum_{k\le n}\varepsilon_k$, where under $P^*$ the $\varepsilon_k=\pm1$ are independent and $+1$ with
probability $p^*$. Taylor expansions give $e^{rh}-e^{-a}=a+(r-\frac{\sigma^2}2)h+O(h^{3/2})$ and $e^a-e^{-a}=2a(1+O(h))$, so
\[
p^*=\frac12+\frac{(r-\frac{\sigma^2}2)\sqrt h}{2\sigma}+O(h).
\]
Hence the mean of $Y_n=\ln\frac{S_T}{S_0}$ is $m_n=na(2p^*-1)=(r-\frac{\sigma^2}2)T+O(\sqrt h)$, and its variance is $s_n^2=4na^2p^*(1-p^*)=\sigma^2T(1-(2p^*-1)^2)\to\sigma^2T$. The
centred and normalised terms $\frac{a(\varepsilon_k-(2p^*-1))}{s_n}$ are bounded by $\frac{2a}{s_n}\to0$, so Lindeberg's condition holds trivially (the truncated sums vanish for
large $n$), and $\frac{Y_n-m_n}{s_n}\to N(0,1)$ in law (Theorem~\ref{thm:lindeberg}). By Slutsky's lemma (Exercise~\ref{ex:slutsky}), $Y_n\to Y$ in law, where $Y$ is
$N(\mu,\sigma^2T)$ with $\mu=(r-\frac{\sigma^2}2)T$.

\emph{The put.} The function $y\mapsto(K-S_0e^y)^+$ is continuous and bounded by $K$, so the binomial put price $e^{-rT}\mathbb E^*[(K-S_0e^{Y_n})^+]$ converges to
$e^{-rT}\mathbb E[(K-S_0e^Y)^+]$. Write $Y=\mu+\sigma\sqrt TZ$ with $Z$ of law $N(0,1)$. The payoff is positive exactly when $Z<-d_-$, and completing the square,
$\mathbb E[e^{\sigma\sqrt TZ}\mathbf 1_{\{Z<-d_-\}}]=e^{\sigma^2T/2}\Phi(-d_--\sigma\sqrt T)$, with $d_-+\sigma\sqrt T=d_+$ and $\mu+\frac{\sigma^2T}2=rT$. So
\[
\begin{aligned}e^{-rT}\mathbb E[(K-S_0e^Y)^+]&=Ke^{-rT}\Phi(-d_-)\eqbr-S_0\Phi(-d_+).\end{aligned}
\]

\emph{The call.} Since $(x-K)^+-(K-x)^+=x-K$ and $\mathbb E^*[S_T]=S_0e^{rT}$ (the discounted price is a $P^*$-martingale), the binomial prices satisfy, for every $n$,
\[
C_n-P_n=e^{-rT}\mathbb E^*[S_T-K]=S_0-Ke^{-rT}.
\]
So $C_n$ converges to $P_{\mathrm{BS}}+S_0-Ke^{-rT}$, which equals $C_{\mathrm{BS}}$ since $1-\Phi(-x)=\Phi(x)$.
\end{proof}

For example, with $S_0=K=100$, $r=5\%$, $\sigma=20\%$ and $T=1$ year, $d_+=0.35$, $d_-=0.15$, and the call is worth $C_{\mathrm{BS}}\approx10.45$ and the put
$P_{\mathrm{BS}}\approx5.57$. The binomial model gives $12.16$ with one period, $10.25$ with $10$, $10.43$ with $100$ and $10.449$ with $1000$: the
convergence is slow and oscillating, at rate about $\frac1n$.

\begin{remark}
Under $P^*$ the binomial price paths converge to $S_t=S_0\exp\bigl((r-\frac{\sigma^2}2)t+\sigma B_t\bigr)$, \emph{geometric Brownian motion}. Black and Scholes, and
Merton, derived the formula in 1973 directly in continuous time, by It\^o's formula and a hedging argument; Bachelier
had used Brownian motion for prices in 1900, without the exponential. The binomial model of 1979 showed that the whole
theory rests on the replication argument of Proposition~\ref{prop:oneperiod} and on the central limit theorem.
\end{remark}

\section{Exercises}

\begin{exercise}[$\star\star$]\label{ex:maxjoint}
\emph{The joint law of the maximum and the endpoint.} Adapt the proof of Theorem~\ref{thm:reflection} to show that for $a>0$ and $b\le a$,
\[
P(M_1\ge a,\ B_1\le b)=P(B_1\ge2a-b).
\]
Deduce the joint density of $(M_1,B_1)$, and show that $M_1-B_1$ also has the law of $|B_1|$.
\end{exercise}

\begin{exercise}[$\star$]\label{ex:bridge}
\emph{The Brownian bridge.} Let $X_t=B_t-tB_1$ for $t\in[0,1]$. Show that $X$ is a centred Gaussian process with $\operatorname{Cov}(X_s,X_t)=\min(s,t)-st$, independent of $B_1$, and
that $B_t=X_t+tZ$ for a variable $Z$ of law $N(0,1)$ independent of $X$. Interpret L\'evy's construction in these terms.
\end{exercise}

\begin{exercise}[$\star\star$]\label{ex:zeros}
\emph{The zeros of Brownian motion.} Let $\mathcal Z=\{t\in[0,1]:B_t=0\}$.
(a) Show that $\mathcal Z$ is closed and has Lebesgue measure $0$ almost surely. (Use Tonelli on $\Omega\times[0,1]$; the map $(t,\omega)\mapsto B_t(\omega)$ is
measurable, being a limit of the step functions $B_{\lceil mt\rceil/m}$.)
(b) Show that almost surely $0$ is not isolated in $\mathcal Z$: $B$ has zeros in $]0,\varepsilon[$ for every $\varepsilon>0$. (Use Theorem~\ref{thm:reflection} and symmetry.)
(c) Show that a non-empty closed subset of $\R$ with no isolated point is uncountable.
\end{exercise}

\begin{exercise}[$\star$]\label{ex:hittingdensity}
Show that the hitting time $T_a$ of Corollary~\ref{cor:hitting} has density $\frac a{\sqrt{2\pi t^3}}e^{-a^2/2t}$ on $]0,\infty[$, and that $T_a$ has the law of $\frac{a^2}{B_1^2}$.
\end{exercise}

\begin{exercise}[$\star$]\label{ex:pvariation}
For the dyadic partitions $\pi_n$ of $[0,1]$ and $p>0$, let $V_p(\pi_n)=\sum_k|B_{k2^{-n}}-B_{(k-1)2^{-n}}|^p$. Compute $\mathbb E[V_p(\pi_n)]$ and $\operatorname{Var}(V_p(\pi_n))$, and show that
$V_p(\pi_n)\to0$ almost surely for $p>2$ and $V_p(\pi_n)\to\infty$ almost surely for $p<2$.
\end{exercise}

\begin{exercise}[$\star$]\label{ex:putcall}
\emph{Put--call parity without a model.} In any market with a bank account at rate $r$, show that if a call and a put with
the same strike $K$ and maturity $T$ trade at prices $C$ and $P$ with $C-P\ne S_0-Ke^{-rT}$, there is an arbitrage. (Compare the payoffs of
``buy the call, sell the put'' and ``buy the stock, borrow $Ke^{-rT}$''.)
\end{exercise}

\begin{exercise}[$\star$]\label{ex:onestep}
A stock worth $100$ will be worth $120$ or $90$ in one year; the bank pays nothing ($r=0$). Price the call with strike $100$ and find the
replicating portfolio. A trader who believes the stock will rise with probability $0.9$ offers to buy the call for $8$. What do you do?
\end{exercise}

\section{Hints}

\begin{hint}{ex:maxjoint}
Reflect the event $\{B_1<b\}$ instead of $\{B_1<a\}$. For the density, differentiate $P(M_1\ge a,B_1\le b)$ in $a$ and $b$. For $M_1-B_1$, integrate the density.
\end{hint}
\begin{hint}{ex:bridge}
Gaussian vectors: compute $\operatorname{Cov}(X_t,B_1)$. At level $n$ of the construction, the displacement of the midpoint is the value of a bridge.
\end{hint}
\begin{hint}{ex:zeros}
(a) $\mathbb E[\lambda(\mathcal Z)]=\int_0^1P(B_t=0)\dd t$. (b) $P(\max_{t\le\varepsilon}B_t>0)=1$ for every $\varepsilon$. (c) Build a binary tree of nested closed intervals, as for
the Cantor set.
\end{hint}
\begin{hint}{ex:hittingdensity}
Differentiate $2(1-\Phi(a/\sqrt t))$. For the second statement, $P(\frac{a^2}{B_1^2}\le t)=P(|B_1|\ge\frac a{\sqrt t})$.
\end{hint}
\begin{hint}{ex:pvariation}
$\mathbb E[V_p]=2^{n(1-p/2)}\mathbb E|Z|^p$ and $\operatorname{Var}(V_p)=2^{n(1-p)}\operatorname{Var}|Z|^p$.
\end{hint}
\begin{hint}{ex:putcall}
$(S_T-K)^+-(K-S_T)^+=S_T-K$ in every scenario.
\end{hint}
\begin{hint}{ex:onestep}
$p^*=\frac{100-90}{120-90}$. Sell the call to the trader and replicate it.
\end{hint}

\section{Solutions}

\begin{solution}{ex:maxjoint}
\idea After the first visit to $a$, the rest of the path is as likely to go up as down by any given amount; reflecting the
final stretch maps ending below $b$ to ending above $2a-b$.
\writeup With the notation of the proof of Theorem~\ref{thm:reflection}, the same argument with the events $\{B_1-B_{T_n}>a-b+\varepsilon\}$ and $\{B_1-B_{T_n}<-(a-b)-\varepsilon\}$ gives
$P(T<1,\ B_1>2a-b+\varepsilon)=P(T<1,\ B_1<b-\varepsilon)$, and letting $\varepsilon\downarrow0$ (the law of $B_1$ has no atoms), $P(T<1,B_1<b)=P(T<1,B_1>2a-b)$. Since $2a-b\ge a$,
$\{B_1>2a-b\}\subseteq\{T<1\}$, so $P(M_1\ge a,\ B_1\le b)=P(B_1\ge2a-b)=1-\Phi(2a-b)$. Differentiating in $b$ and then in $a$ (with $\varphi=\Phi'$, $\varphi'(x)=-x\varphi(x)$), the
joint density on $\{a>0,\ b<a\}$ is
\[
f(a,b)=2(2a-b)\varphi(2a-b).
\]
For $M_1-B_1$: $P(M_1-B_1\ge c)=\int\!\!\int_{a-b\ge c}2(2a-b)\varphi(2a-b)\dd b\dd a$. With $x=a-b\ge c$ and $a>0$, $2a-b=a+x$, and $\int_0^\infty2(a+x)\varphi(a+x)\dd a=2\varphi(x)$, so
$P(M_1-B_1\ge c)=\int_c^\infty2\varphi(x)\dd x=P(|B_1|\ge c)$. (This is a shadow of L\'evy's theorem that $M-B$ is, as a process, a reflected Brownian motion.)
\end{solution}

\begin{solution}{ex:bridge}
\idea A Gaussian vector splits into independent pieces as soon as the covariances vanish: subtracting the right multiple
of $B_1$ kills the correlation.
\writeup The values of $X$ and $B_1$ at finitely many times form a Gaussian vector (Theorem~\ref{thm:gaussvec}(d)), centred, with
$\operatorname{Cov}(X_s,X_t)=\min(s,t)-st-ts+st=\min(s,t)-st$ and $\operatorname{Cov}(X_t,B_1)=t-t=0$. So $(X_{t_1},\dots,X_{t_k})$ is independent of $B_1$ for all times (Theorem~\ref{thm:gaussvec}(a) and
Theorem~\ref{thm:inversion}(a), as in Proposition~\ref{prop:scaling}), and $B_t=X_t+tB_1$ with $Z=B_1$. The bridge is Brownian motion conditioned to return to $0$ at time $1$.
In L\'evy's construction, between two consecutive points of level $n-1$ the path is a scaled bridge added to the linear interpolation; its value at the
midpoint is $N(0,\frac14\cdot2^{-n+1})=N(0,2^{-n-1})$, of standard deviation $2^{-(n+1)/2}$: exactly the displacement used.
\end{solution}

\begin{solution}{ex:zeros}
\idea Zero has probability zero at each fixed time, hence a zero-measure set of zeros; but near time $0$ the path
changes sign infinitely often.
\writeup (a) $\mathcal Z$ is the preimage of $\{0\}$ by a continuous function, so it is closed. By Tonelli (A2), $\mathbb E[\lambda(\mathcal Z)]=\int_0^1P(B_t=0)\dd t=0$, so $\lambda(\mathcal Z)=0$ almost surely.
(b) By Theorem~\ref{thm:reflection} and scaling, $P(\max_{t\le\varepsilon}B_t>0)\ge P(\max_{t\le\varepsilon}B_t\ge\delta)=P(|B_\varepsilon|\ge\delta)\to1$ as $\delta\to0$; so the event $E_\varepsilon=\{\max_{t\le\varepsilon}B_t>0\}$ has
probability $1$, and so has its analogue for $-B$. Intersecting over $\varepsilon=\frac1j$, almost surely $B$ takes positive and negative values in every
interval $]0,\varepsilon[$, hence vanishes somewhere in it by the intermediate value theorem.
(c) Let $F$ be non-empty, closed, without isolated points. Pick two distinct points of $F$ and disjoint closed intervals of length
at most $1$ around them; each interval meets $F$ in a point that is not isolated, so each contains two points of $F$ again, around which
we put two disjoint closed intervals of length at most $\frac12$ inside it, and so on. Each infinite sequence of left/right choices gives
a decreasing sequence of closed intervals with lengths tending to $0$, each meeting the closed set $F$; their intersection is a point of
$F$ (nested closed intervals with lengths tending to $0$ meet in one point, by completeness of $\R$). Different sequences give different points, so $F$
has at least as many points as $\{0,1\}^{\N}$, which is uncountable (Cantor's diagonal argument). Every point of $\mathcal Z$ is in fact a non-isolated point; this uses the strong Markov
property at the first zero after each rational time, proved in Tome~11, and with (c) it shows that $\mathcal Z$ is uncountable.
\end{solution}

\begin{solution}{ex:hittingdensity}
\idea The law of the hitting time is read off from the law of the maximum, by the scaling relation between space and
time.
\writeup By Corollary~\ref{cor:hitting}, $P(T_a\le t)=2(1-\Phi(at^{-1/2}))$. Differentiating, the density is $2\varphi(at^{-1/2})\cdot\frac a2t^{-3/2}=\frac a{\sqrt{2\pi t^3}}e^{-a^2/2t}$. And
$P(\frac{a^2}{B_1^2}\le t)=P(|B_1|\ge at^{-1/2})=2(1-\Phi(at^{-1/2}))$: the same distribution function. The density decays like $t^{-3/2}$, which again shows $\mathbb E[T_a]=\infty$.
\end{solution}

\begin{solution}{ex:pvariation}
\idea Each of the $2^n$ increments is $2^{-n/2}$ times a standard Gaussian, so the $p$-th powers sum to $2^{n(1-p/2)}$ times an average,
which the law of large numbers controls.
\writeup The increments are $2^{-n/2}Z_k$ with $Z_k$ independent $N(0,1)$, so $V_p(\pi_n)=2^{-np/2}\sum_{k\le2^n}|Z_k|^p$, of mean $c_p2^{n(1-p/2)}$ with $c_p=\mathbb E|Z|^p\in\,]0,\infty[$, and
variance $v_p2^{n(1-p)}$ with $v_p=\operatorname{Var}|Z|^p$. By Chebyshev, $P\bigl(|V_p-\mathbb E V_p|>\frac12\mathbb E V_p\bigr)\le\frac{4v_p2^{n(1-p)}}{c_p^22^{n(2-p)}}=\frac{4v_p}{c_p^2}2^{-n}$, which is summable; by
Borel--Cantelli, almost surely $\frac12\mathbb E V_p\le V_p\le\frac32\mathbb E V_p$ for large $n$. Since $\mathbb E V_p\to0$ for $p>2$ and $\to\infty$ for $p<2$, the claims follow. The exponent $2$ is
the dividing line, in agreement with Theorem~\ref{thm:qv}.
\end{solution}

\begin{solution}{ex:putcall}
\idea Two portfolios with the same payoff in every scenario must cost the same today; otherwise buy the cheap one, sell
the dear one, and pocket the difference.
\writeup Portfolio A: buy the call and sell the put; at $T$ it is worth $(S_T-K)^+-(K-S_T)^+=S_T-K$. Portfolio B: buy one share and borrow $Ke^{-rT}$; at $T$
it is worth $S_T-K$ too. Their costs are $C-P$ and $S_0-Ke^{-rT}$. If $C-P>S_0-Ke^{-rT}$, sell A and buy B: you receive the difference today and owe
nothing at $T$, whatever happens. If $C-P<S_0-Ke^{-rT}$, do the opposite. No model of the stock is needed.
\end{solution}

\begin{solution}{ex:onestep}
\idea The option is a combination of stock and cash; its price is the price of that combination, whatever anybody
believes about the future.
\writeup Here $u=1.2$, $d=0.9$, $e^{rh}=1$, so $p^*=\frac{1-0.9}{1.2-0.9}=\frac13$. The call pays $20$ or $0$, so its price is $\frac13\cdot20=\frac{20}3\approx6.67$. Replication: $\delta=\frac{20-0}{120-90}=\frac23$
share, financed by borrowing $\frac23\cdot90=60$: the portfolio costs $\frac{200}3-60=\frac{20}3$ and is worth $80-60=20$ or $60-60=0$. The trader's beliefs are irrelevant:
sell him the call for $8$, buy $\frac23$ of a share and borrow $60$; you keep $8-\frac{20}3=\frac43$ for sure, and the portfolio pays exactly what you owe him.
\end{solution}

\section*{Sources and further reading}

\begin{itemize}
\item P.~M\"orters and Y.~Peres, \emph{Brownian Motion}, Cambridge University Press, 2010. L\'evy's construction, scaling,
nowhere differentiability, the reflection principle and the zero set, all in Chapters~1 and~2.
\item I.~Karatzas and S.~Shreve, \emph{Brownian Motion and Stochastic Calculus}, 2nd ed., Springer, 1991.
\item P.~L\'evy, \emph{Processus stochastiques et mouvement brownien}, Gauthier-Villars, 1948.
\item N.~Wiener, \emph{Differential space}, Journal of Mathematics and Physics 2 (1923), 131--174.
\item R.~Paley, N.~Wiener and A.~Zygmund, \emph{Notes on random functions}, Mathematische Zeitschrift 37 (1933), 647--668.
\item A.~Dvoretzky, P.~Erd\H os and S.~Kakutani, \emph{Nonincrease everywhere of the Brownian motion process}, Proceedings of
the Fourth Berkeley Symposium, Vol.~2, 1961, 103--116.
\item A.~Einstein, \emph{\"Uber die von der molekularkinetischen Theorie der W\"arme geforderte Bewegung von in ruhenden
Fl\"ussigkeiten suspendierten Teilchen}, Annalen der Physik 17 (1905), 549--560.
\item L.~Bachelier, \emph{Th\'eorie de la sp\'eculation}, Annales scientifiques de l'\'Ecole Normale Sup\'erieure 17 (1900), 21--86.
\item J.~Cox, S.~Ross and M.~Rubinstein, \emph{Option pricing: a simplified approach}, Journal of Financial Economics 7
(1979), 229--263.
\item F.~Black and M.~Scholes, \emph{The pricing of options and corporate liabilities}, Journal of Political Economy 81
(1973), 637--654.
\item S.~Shreve, \emph{Stochastic Calculus for Finance I: The Binomial Asset Pricing Model}, Springer, 2004.
\end{itemize}
